\documentclass[10pt,a4paper]{amsart}
\usepackage[margin=19mm]{geometry}
\usepackage{amsmath,amssymb,mathtools}
\usepackage[scr=rsfs,cal=euler]{mathalfa}
\usepackage{xfrac}
\usepackage[unicode,hidelinks,bookmarks=false]{hyperref}
\newcommand{\ie}{i.e.\ }
\newcommand{\cf}{cf.\ }
\newcommand{\resp}{respectively\ }
\newcommand{\Subsection}{\subsection}
\providecommand{\nobreakdash}{\nobreak\hbox{-}}
\setlength{\emergencystretch}{2em}
\multlinegap0pt
\usepackage{amssymb}
\newtheorem{proposition}{Proposition}
\title{Attractor Flow Versus Hesse Flow in Wall-Crossing Structures}
\author{Qiang Wang}
\date{Source: arXiv:2607.05433v1 (CC BY 4.0)}
\begin{document}
\maketitle

\noindent\emph{Source and licence.} Attractor Flow Versus Hesse Flow in Wall-Crossing Structures. Version arXiv:2607.05433v1 (CC BY 4.0), distributed by arXiv under the Creative Commons Attribution 4.0 International licence, http://creativecommons.org/licenses/by/4.0/, as recorded in the arXiv licence metadata for this exact version and shown on the abstract page https://arxiv.org/abs/2607.05433v1. Full licence notice: \texttt{licenses/arxiv-2607.05433-CC-BY-4.0.txt}.\par\smallskip

\noindent\textit{Editorial scope.} Counted unit: the article's proposition pro:4.3 (source0.tex lines 381--386). The article's complete original proof of it is delivered with it (source0.tex lines 388--401, one \texttt{proof} environment of 14 lines). No mathematics is written, repaired or re-derived: the statement and the proof are the article's own lines, in the article's own notation, and no retained line is substituted or re-flowed.  No further source-local context is needed: every cross-reference of every retained line resolves inside the two delivered spans.  Nothing was omitted from any retained span, so the package carries no omission record.  No retained line contains a citation, so the wrapper carries no bibliography and no bibliography entry.  No retained line contains a figure environment, an image-inclusion command or drawing code, so no image is delivered and the package declares no asset dependency.  Editorial formatting only, in addition: an AMS \texttt{amsart} preamble that restates the article's own declarations and macros used by the retained lines, namely the environments \texttt{proposition} on separate counters, together with the standard packages those lines need.  The retained environments print in this document as Proposition 1 in order of appearance, whereas the article prints the same items with the numbering of its own full document.  The complete native source of the article as distributed in the arXiv e-print for this exact version is bundled unchanged as \texttt{sources/204542/source0.tex}.
\begin{proposition}\label{pro:4.3}
The affine variables $x_{l}$ is conjugate to $y^{l}$, while $y_{l}$ is conjugate to $x^{l}$ in the sense that there exists some holomorphic function $\mathscr{F}$ such that
\begin{equation}\label{eq:18}
          x_{l}=\frac{\partial\mathscr{F}}{\partial y^{l}},\qquad y_{l}=\frac{\partial\mathscr{F}}{\partial x^{l}}
\end{equation}
\end{proposition}

\begin{proof} We compute as follows:
\[x_{l}+iy_{l}=e^{-i\theta}a_{D,l}=e^{-i\theta}\frac{\partial\mathcal{F}}{\partial a^{l}}=e^{-i\theta}\left(\sum_{k}\frac{\partial\mathcal{F}}{\partial x^{k}}\frac{\partial x^{k}}{\partial a^{l}}+\sum_{k}\frac{\partial\mathcal{F}}{\partial y^{k}}\frac{\partial y^{k}}{\partial a^{l}}\right)=
\]

\[=e^{-i\theta}\left(e^{-i\theta}\frac{\partial\mathcal{F}}{\partial x^{l}}-i\,e^{-i\theta}\frac{\partial\mathcal{F}}{\partial y^{l}}\right)=\frac{\partial}{\partial x^{l}}\left(e^{-2i\theta}\mathcal{F}\right)-i\,\frac{\partial}{\partial y^{l}}\left(e^{-2i\theta}\mathcal{F}\right)=\]

\[=\frac{\partial}{\partial x^{l}}\left(Re\left(e^{-2i\theta}\mathcal{F}\right)+i\,Im\left(e^{-2i\theta}\mathcal{F}\right)\right)-i\,\frac{\partial}{\partial y^{l}}\left(Re\left(e^{-2i\theta}\mathcal{F}\right)+i\,Im\left(e^{-2i\theta}\mathcal{F}\right)\right)\]

which by applying the Cauchy-Riemann equation for the holomorphic function $e^{-2i\theta}\mathcal{F}$, becomes

\[2\left(\frac{\partial Im\left(e^{-2i\theta}\mathcal{F}\right)}{\partial y^{l}}+i\,\frac{\partial Im\left(e^{-2i\theta}\mathcal{F}\right)}{\partial x^{l}}\right)
\]
By defining $\mathscr{F}:=2Im\left(e^{-2i\theta}\mathcal{F}\right)$, the relations (\ref{eq:18}) then follows from the above computation.
\end{proof}

\end{document}
